% ============================================================
%  06_kiem_thu_tu_dong.tex  --  Kiem thu tu dong
% ============================================================

\section{Kết quả kiểm thử tự động}
\label{sec:tudong}

Bộ kiểm thử tự động được tổ chức trong file
\textit{tests/test\_memory\_safety.py}. Script này tích hợp cả hai luồng
kiểm thử: phát hiện lỗi runtime bằng ASan và xác nhận tính nhất quán của
mô hình đánh giá rủi ro.

% ────────────────────────────────────────────────────────────
\subsection{Cấu trúc bộ kiểm thử}

\subsubsection{Khởi tạo -- \textit{setup\_module()}}

Trước khi chạy bất kỳ bài kiểm thử nào, \textit{setup\_module()} biên dịch
cả hai binary bằng GCC với cờ \textit{-fsanitize=address}. Điều này đảm bảo
mọi kết quả kiểm thử đều dựa trên binary mới nhất từ source hiện tại.

\begin{lstlisting}[language=Python, caption={Hàm biên dịch và chạy binary}]
def _build(src_name: str, bin_name: str):
    subprocess.run(
        ["gcc", "-g", "-fsanitize=address", "-Wall", "-Wextra",
         "-o", bin_name, src_name],
        cwd=C_DIR, check=True,
    )

def _run(bin_name: str, args: list[str]):
    return subprocess.run(
        [os.path.join(C_DIR, bin_name)] + args,
        cwd=C_DIR, capture_output=True, text=True,
    )
\end{lstlisting}

\subsubsection{Nhóm 1: Xác nhận lỗi trên bản vulnerable (3 bài)}

\begin{description}
    \item[\textit{test\_vulnerable\_overflow\_detected}]
          Chạy \textit{vulnerable\_inventory overflow "A"*39}. Kiểm tra:
          \textit{returncode != 0} (ASan buộc dừng) và chuỗi
          \textit{"heap-buffer-overflow"} xuất hiện trong \textit{stderr}.

    \item[\textit{test\_vulnerable\_uaf\_detected}]
          Chạy \textit{vulnerable\_inventory uaf}. Kiểm tra:
          \textit{returncode != 0} và \textit{"heap-use-after-free"} trong
          \textit{stderr}.

    \item[\textit{test\_vulnerable\_leak\_detected}]
          Chạy \textit{vulnerable\_inventory leak} với
          \textit{ASAN\_OPTIONS=detect\_leaks=1}. Kiểm tra: chuỗi
          \textit{"leaked"} xuất hiện trong \textit{stderr} (không yêu cầu
          exit code khác 0 vì LeakSanitizer có thể không dừng chương trình
          trên mọi platform).
\end{description}

\subsubsection{Nhóm 2: Xác nhận bản secure sạch (3 bài)}

\begin{description}
    \item[\textit{test\_secure\_overflow\_clean}]
          Chạy \textit{secure\_inventory overflow "A"*39}. Kiểm tra:
          \textit{returncode == 0} và \textit{"ERROR"} \emph{không} xuất
          hiện trong \textit{stderr}.

    \item[\textit{test\_secure\_uaf\_clean}]
          Chạy \textit{secure\_inventory uaf}. Kiểm tra:
          \textit{returncode == 0} và \textit{"ERROR"} không trong
          \textit{stderr}.

    \item[\textit{test\_secure\_leak\_clean}]
          Chạy \textit{secure\_inventory leak} với
          \textit{ASAN\_OPTIONS=detect\_leaks=1}. Kiểm tra:
          \textit{returncode == 0} và \textit{"leaked"} không trong
          \textit{stderr}.
\end{description}

\subsubsection{Nhóm 3: Kiểm tra mô hình rủi ro (2 bài)}

\begin{description}
    \item[\textit{test\_risk\_assessment\_matches\_baseline}]
          Gọi \textit{assess\_all()} để tính điểm rủi ro từ dataset, đọc
          \textit{expected\_risk\_baseline.json}, so sánh từng cặp kết quả:
          \textit{id}, \textit{severity} phải khớp chính xác; \textit{base\_score}
          được phép sai lệch tối đa 0.05 điểm.

    \item[\textit{test\_critical\_vuln\_is\_the\_buffer\_overflow}]
          Lấy phần tử đầu (điểm cao nhất) từ kết quả \textit{assess\_all()}.
          Xác nhận \textit{severity == "CRITICAL"} và \textit{cwe == "CWE-121"}.
\end{description}

% ────────────────────────────────────────────────────────────
\subsection{Log kiểm thử}

\begin{lstlisting}[language=bash, caption={Log make test -- 8/8 PASS}]
$ make test
gcc -g -Wall -Wextra -fsanitize=address \
    -o c/vulnerable_inventory c/vulnerable_inventory.c
gcc -g -Wall -Wextra -fsanitize=address \
    -o c/secure_inventory c/secure_inventory.c
python3 tests/test_memory_safety.py

[PASS] test_vulnerable_overflow_detected
[PASS] test_vulnerable_uaf_detected
[PASS] test_vulnerable_leak_detected
[PASS] test_secure_overflow_clean
[PASS] test_secure_uaf_clean
[PASS] test_secure_leak_clean
[PASS] test_risk_assessment_matches_baseline
[PASS] test_critical_vuln_is_the_buffer_overflow

8/8 test da PASS
\end{lstlisting}

% ────────────────────────────────────────────────────────────
\subsection{Ảnh log kiểm thử thực tế}

\begin{figure}[H]
    \centering
    \includegraphics[width=\textwidth]{log_auto_test}
    \caption{Log kiểm thử tự động -- toàn bộ 8/8 PASS}
\end{figure}

% ────────────────────────────────────────────────────────────
\subsection{Phân tích kết quả từng nhóm}

\subsubsection{Nhóm 1: Xác nhận phát hiện lỗi}

Ba bài kiểm thử nhóm 1 đều PASS, xác nhận:
\begin{itemize}
    \item ASan phát hiện đủ 3/3 lỗi trên bản vulnerable với các kịch bản được
          thiết kế.
    \item Chuỗi nhận dạng lỗi trong stderr khớp chính xác với loại lỗi được
          khai thác:
          \begin{itemize}
              \item \textit{"heap-buffer-overflow"} cho CWE-121
              \item \textit{"heap-use-after-free"} cho CWE-416
              \item \textit{"leaked"} cho CWE-401
          \end{itemize}
    \item Tiến trình kết thúc với exit code khác 0 (trường hợp overflow và UAF),
          xác nhận ASan buộc dừng chương trình ngay khi phát hiện lỗi.
\end{itemize}

\subsubsection{Nhóm 2: Xác nhận bản secure sạch}

Ba bài kiểm thử nhóm 2 đều PASS, xác nhận:
\begin{itemize}
    \item Bản secure 0/3 lỗi với cùng ba kịch bản.
    \item Exit code 0 trong cả ba trường hợp -- chương trình hoàn thành bình
          thường.
    \item Không có thông báo ERROR trong stderr -- không có vi phạm bộ nhớ nào
          được ASan phát hiện.
    \item Riêng kịch bản overflow: chuỗi bị cắt bớt an toàn xuống tối đa 31
          ký tự, không gây tràn.
\end{itemize}

\subsubsection{Nhóm 3: Tính nhất quán của mô hình rủi ro}

Hai bài kiểm thử nhóm 3 đều PASS, xác nhận:
\begin{itemize}
    \item Kết quả tính toán của \textit{risk\_assessor.py} khớp với baseline
          (\textit{expected\_risk\_baseline.json}) với sai lệch $\leq 0.05$
          điểm cho tất cả 5 lỗ hổng.
    \item Phần tử điểm cao nhất trong danh sách kết quả là VULN-001 (CWE-121)
          với mức CRITICAL -- đúng như dự kiến từ phân tích lý thuyết.
    \item Thứ tự sắp xếp theo điểm giảm dần nhất quán:
          VULN-001 = VULN-005 (9.8) $>$ VULN-002 (6.9) $>$ VULN-003 (5.4)
          $>$ VULN-004 (2.5).
\end{itemize}

% ────────────────────────────────────────────────────────────
\subsection{Nhận xét về thiết kế bộ kiểm thử}

Bộ kiểm thử có một số điểm thiết kế đáng chú ý:

\begin{itemize}
    \item Tự biên dịch trước khi chạy: \textit{setup\_module()} đảm bảo binary
          luôn đồng bộ với source, tránh kết quả sai do dùng binary cũ.
    \item Tách biệt môi trường cho kịch bản leak: Kịch bản leak dùng
          \textit{env = dict(os.environ, ASAN\_OPTIONS="detect\_leaks=1")} thay
          vì set biến môi trường toàn cục, tránh ảnh hưởng đến các kịch bản khác.
    \item Kiểm tra cả ASan lẫn risk model: Bộ test không chỉ kiểm tra runtime
          mà còn xác nhận tính nhất quán của mô hình đánh giá rủi ro, tạo thành
          một vòng xác minh khép kín.
    \item Hạn chế: Bộ test chỉ kiểm tra các kịch bản được lập trình sẵn; không
          kiểm tra tràn trong struct (đây là giới hạn của ASan đã phân tích ở
          Phần~\ref{sec:hopden}). Các edge case khác (payload rỗng, ký tự đặc
          biệt) không được bao phủ.
\end{itemize}
